\subsection{Square-integrable representations}

In this issue, the Hilbert spaces considered are complex. It is assumed that G is unimodular, and its center is denoted by C. For every closed subgroup Z of C, we denote by \(\beta_{Z}\) a Haar measure on Z, and we equip G/Z with the Haar measure \(\nu_{Z} = \mu/\beta_{Z}\) (INT, VII, p. 44, § 2, no. 2, def. 1).

Let \(\pi\) an irreducible unitary representation of G in a Hilbert space E and \(\chi \in \widehat{\mathbf{C}}\) its central character (def. 6 of V, p. 390). For all \(x\) and \(y\) in E, we denote \(f_{x,y}\) the matrix coefficient \(g \mapsto \langle x \mid \pi(g)y \rangle\) ; it is a continuous function on G with complex values.

Let \(x\) and \(y\) in E. For \(z \in \mathrm{C}\) and \(g \in \mathrm{G}\) , we have

\[
f _ {x, y} (z g) = \langle x \mid \pi (z g) y \rangle = \langle x \mid \chi (z) \pi (g) y \rangle = \chi (z) f _ {x, y} (g),\tag{9}
\]

hence \(f_{x,y}\) belongs to the space \(\mathcal{F}_{\chi}(\mathrm{G})\) (V, p. 408). Moreover, for every \((g_1, h_1) \in \mathrm{G} \times \mathrm{G}\) and every \(g \in \mathrm{G}\) , we have

\[
f _ {\pi (g _ {1}) x, \pi (h _ {1}) y} (g) = \langle \pi (g _ {1}) x | \pi (g) \pi (h _ {1}) y \rangle = f _ {x, y} (g _ {1} ^ {- 1} g h _ {1}).\tag{10}
\]

Relation (9) justifies the following definition.

DEFINITION 3. --- Let \(\pi\) an irreducible unitary representation of G in a Hilbert space E. We say that \(\pi\) is square-integrable modulo the center if the function on G/C deduced from the function \(|f_{x,y}|\) by passage to the quotient belongs to \(\mathcal{L}^2 (\mathrm{G / C})\) for all \((x,y)\in \mathrm{E}\times \mathrm{E}\) .

This condition does not depend on the choice of a Haar measure on G/C.

If the matrix coefficients of \(\pi\) belong to \(\mathcal{L}^2 (\mathrm{G})\) , one says that \(\pi\) is square-integrable; the existence of an irreducible unitary representation of G that is square-integrable implies that the center of G is compact (exercise 5 of V, p. 487).

There exist groups G which admit no square-integrable representation, even modulo the center (cf. exercise 32 of V, p. 516).

PROPOSITION 6. --- Let \(\pi\) an irreducible unitary representation of G in a Hilbert space E. Let \(\chi\) the central character of \(\pi\) . Then \(\pi\) is square-integrable modulo the center if and only if there exist nonzero elements \(x_0\) and \(y_0\) of E such that the function on G/C induced by \(|f_{x_0,y_0}|\) by passage to the quotient belongs to \(\mathcal{L}^2 (\mathrm{G / C})\) .

Suppose there exist elements \(x_{0}\) and \(y_{0}\) nonzero elements of E such that \(|f_{x_{0},y_{0}}| \in \mathcal{L}^{2}(\mathrm{G/C})\) . It suffices to show that \(\pi\) is then square-integrable modulo the center.

 Let F be the set of elements \(x \in E\) such that \(|f_{x,y_{0}}|\) belongs to \(\mathcal{L}^{2}(G/C)\) . This is a vector subspace of E; it contains \(x_{0}\) , hence is nonzero. Relation (10) above implies that F is stable under \(\pi\) ; since the representation \(\pi\) is irreducible, the space F is therefore dense in E.

 Let \(x \in \mathrm{F}\) . Since \(f_{x,y_0}\) belongs to \(\mathcal{F}_{\chi}(\mathrm{G})\) and is \(\mu\) -measurable, and \(|f_{x,y_0}| \in \mathcal{L}^2(\mathrm{G/C})\) , we have \(f_{x,y_0} \in \mathcal{L}_{\chi}^2(\mathrm{G})\) (prop. 5 of V, p. 413 applied to Z = C). Let \(u\) denote the partial operator from E to \(\mathrm{L}_{\chi}^{2}(\mathrm{G})\) whose domain is F and which associates to \(x \in \mathrm{F}\) the class of \(f_{x,y_0}\) in \(\mathrm{L}_{\chi}^{2}(\mathrm{G})\) .

 Let us show that the partial operator \(u\) is closed. Let \((x_{n}, u(x_{n}))_{n \in \mathbf{N}}\) a sequence of elements of the graph of \(u\) which converges in \(\mathrm{E} \times \mathrm{L}_{\chi}^{2}(\mathrm{G})\) . Let \(x\) the limit of \((x_{n})\) and let \(f \in \mathcal{L}_{\chi}^{2}(\mathrm{G})\) a function whose class is the limit of the sequence \((u(x_{n}))\) .

The function \(u(x_{n})\) is the class of the matrix coefficient \(f_{x_{n},y_{0}}\) . For every \(g \in G\) , we have

\[
f _ {x _ {n}, y _ {0}} (g) = \langle x _ {n} | \pi (g) y _ {0} \rangle \rightarrow \langle x | \pi (g) y _ {0} \rangle = f _ {x, y _ {0}} (g)
\]

as \(n \to +\infty\) . We therefore have \(f_{x,y_{0}} \in \mathcal{L}_{\chi}^{2}(\mathrm{G})\) and \(f = f_{x_{0},y}\) almost everywhere modulo C (cor. of prop. 3 of V, p. 409); this means that \(u(x)\) is the class of f in \(\mathrm{L}_{\chi}^{2}(\mathrm{G})\) . We have therefore proved that u is closed.

The domain F of u is stable under \(\pi\) , and relation (10) shows that u satisfies \(u \circ \pi(g) = \gamma_{\mathrm{G},\chi}(g) \circ u\) for all \(g \in G\) . Consequently we also have the equality \(u^{*} \circ \gamma_{\mathrm{G},\chi}(g) = \pi(g) \circ u^{*}\) for all \(g \in G\) (lemma 6 of V, p. 381), whence \((u^{*} \circ u) \circ \pi(g) = \pi(g) \circ (u^{*} \circ u)\) . Now the partial operator \(u^{*} \circ u\) is self-adjoint (prop. 12 of IV, p. 241), and in particular closed (prop. 7 of IV, p. 236), so cor. 1 of V, p. 387 implies that the domain of \(u^{*} \circ u\) is equal to E. A fortiori, we have F = E, that is, the function \(|f_{x,y_{0}}|\) belongs to \(\mathcal{L}^{2}(\mathrm{G}/\mathrm{C})\) for all \(x \in E\) .

Let \((x,y)\in\mathrm{E}\times\mathrm{E}\) . We have \(f_{x,y_{0}}\in\mathcal{L}_{\chi}^{2}(\mathrm{G})\) . One proves mutatis mutandis, using the representation \(\delta_{G,\chi}\) instead of \(\gamma_{G,\chi}\) , that the set of \(y\in E\) such that the function \(|f_{x,y}|\) belongs to \(\mathcal{L}^{2}(\mathrm{G}/\mathrm{C})\) is equal to E. The proposition follows.

In the rest of this section, we fix a closed subgroup Z of C such that C/Z is compact.

Lemma 9. --- Let \(\pi\) an irreducible unitary representation of G in a Hilbert space E which is square-integrable modulo the center. Let \(\chi\) the restriction to Z of the central character of \(\pi\) . For all \(x\) and \(y\) in E, we have \(f_{x,y} \in \mathcal{L}_{\chi}^{2}(\mathrm{G})\) .

 The function \(f_{x,y}\) is continuous hence \(\mu\) -measurable. We have \(f_{x,y} \in \mathcal{F}_{\chi}(\mathrm{G})\) by formula (9), p. 420. Moreover, by INT, VII, p. 64, § 2, no. 8, cor. 1, c), we have \(\mathrm{N}_{2}(f_{x,y}) < +\infty\) since C/Z is compact and \(|f_{x,y}| \in \mathcal{L}^{2}(\mathrm{G}/\mathrm{C})\) . The assertion therefore follows from Proposition 5 of V, p. 413.

PROPOSITION 7. --- Let π be an irreducible unitary representation of G in a Hilbert space E which is square-integrable modulo the center. Let χ be the restriction to Z of the central character of π.

There exists a real number c \textgreater{} 0 and a unique \((\mathbf{G} \times \mathbf{G})\) -isometric morphism w of the unitary representation \(\overline{\pi} \boxtimes \pi\) in \(\mathrm{L}_{\chi}^{2}(\mathrm{G})\) such that, for all \((x, y) \in \overline{\mathrm{E}} \times \mathrm{E}\) , the element \(w(x \otimes y)\) is the class in \(\mathrm{L}_{\chi}^{2}(\mathrm{G})\) of the function \(c^{1/2} f_{x,y}\) .

 For every \((x,y)\in\mathrm{E}\times\mathrm{E}\) , we have \(f_{x,y}\in\mathcal{L}_{\chi}^{2}(\mathrm{G})\) (Lemma 9). Let v denote the unique linear map from \(\overline{E}\otimes E\) in \(\mathrm{L}_{\chi}^{2}(\mathrm{G})\) such that \(v(x\otimes y)\) is the class of \(f_{x,y}\) for all \((x,y)\in\overline{\mathrm{E}}\times\mathrm{E}\) .

We prove below the following lemma.

Lemma 10. --- There exists a real number c \textgreater{} 0 such that the linear map \(w = c^{1/2}v\) is isometric.

This lemma being assumed valid, note that formula (10), p. 420, can be written as

\[
v (\overline {{\pi}} (g _ {1}) x \otimes \pi (h _ {1}) y) = \varrho_ {\mathrm{G}, \chi} (g _ {1}, h _ {1}) v (x \otimes y)\tag{11}
\]

 for all \((g_{1}, h_{1}) \in \mathbf{G} \times \mathbf{G}\) and every \((x, y) \in \overline{\mathbf{E}} \times \mathbf{E}\) . The isometric linear map w from \(\overline{\mathbf{E}} \otimes \mathbf{E}\) in \(\mathrm{L}_{\chi}^{2}(\mathrm{G})\) admits a continuous extension, still denoted w, to \(\overline{E} \widehat{\otimes}_{2} E\) . By continuity and linearity, formula (11) implies that w is a \((\mathrm{G} \times \mathrm{G})\) -morphism of \(\overline{\pi} \boxtimes \pi\) in \(\mathrm{L}_{\chi}^{2}(\mathrm{G})\) , which concludes the proof of the proposition.

 Let us prove the lemma. For all \(x \in E\) , we denote \(u_{x}\) the linear map \(y \mapsto v(x \otimes y) = f_{x,y}\) from E into \(\mathrm{L}_{\chi}^{2}(\mathrm{G})\) . We have \(u_{x} \in \mathrm{Hom}_{\mathrm{G}}(\pi, \gamma_{\mathrm{G},\chi})\) . By formula (10), according to Corollary 5 of V, p. 388, there exists a real number \(\lambda_{x} \geqslant 0\) such that \(\lambda_{x}u_{x}\) is isometric.

Let \(x\) and \(y\) in E. We have

\[
\| f _ {x, y} \| ^ {2} = \int_ {\mathrm{G/Z}} \overline {{f}} _ {x, y} f _ {x, y} d \nu_ {\mathrm{Z}} = \int_ {\mathrm{G/Z}} \check {\overline {{f}}} _ {x, y} \check {f} _ {x, y} d \nu_ {\mathrm{Z}}
\]

since G/Z is unimodular (Lemma 7 of V, p. 417). Noting that \(\check{f}_{x,y} = \overline{f}_{y,x}\) , we obtain

\[
\lambda_ {x} \| y \| ^ {2} = \| f _ {x, y} \| ^ {2} = \int_ {\mathrm{G/Z}} f _ {y, x} \overline {{f}} _ {y, x} d \nu_ {\mathrm{Z}} = \| f _ {y, x} \| ^ {2} = \lambda_ {y} \| x \| ^ {2}.
\]

This means that the positive real number \(\lambda_{x}/\|x\|^{2}\) is independent of the choice of the nonzero element x of E. It is strictly positive since, for every nonzero x, the function \(f_{x,x}\) is continuous and takes the value \(\|x\|^{2}>0\) at e, whence \(\|u_{x}(x)\|=\|f_{x,x}\|>0\) . Denote by \(c^{-1}\) this real number.

For every \((x,y)\in\mathrm{E}\times\mathrm{E}\) , it follows that

\[
\| v (x \otimes y) \| ^ {2} = \| f _ {x, y} \| ^ {2} = \lambda_ {x} \| y \| ^ {2} = c ^ {- 1} \| x \| ^ {2} \| y \| ^ {2} = c ^ {- 1} \| x \otimes y \| ^ {2}.
\]

Using EVT, V, p. 29, cor. 1, we deduce that the linear map \(w = c^{1/2}v\) of \(\overline{\mathrm{E}}\widehat{\otimes}_2\mathrm{E}\) in \(\mathrm{L}_{\chi}^{2}(\mathrm{G})\) is isometric, as desired.

 COROLLARY. --- Let \(\pi\) an irreducible unitary representation of G and let \(\chi\) the restriction to Z of its central character. The representation \(\pi\) is square-integrable modulo the center if and only if it is isomorphic to a subrepresentation of the representation \(\gamma_{\mathrm{G},\overline{\chi}}\) from G into \(\mathrm{L}_{\overline{\chi}}^{2}(\mathrm{G})\) (resp. of the representation \(\delta_{\mathrm{G},\chi}\) from G into \(\mathrm{L}_{\chi}^{2}(\mathrm{G})\) ).

Let us prove the assertion concerning \(\gamma_{\mathrm{G},\overline{\chi}}\) , the second being proved in a similar manner.

 Suppose first that there exists a subrepresentation E of the representation \(\gamma_{\mathrm{G},\overline{\chi}}\) isomorphic to \(\pi\) . Since the space E is not zero, there exists a function \(f_1 \in \mathcal{K}_{\overline{\chi}}(\mathrm{G})\) whose class \(\widetilde{f}_1\) that is not orthogonal to E. We have \(f_1 \in \mathcal{L}_{\overline{\chi}}^1 (\mathrm{G})\) . Denote by \(\widetilde{f}_{1,\mathrm{E}}\) the orthogonal projection of \(\widetilde{f}_1\) on E; it is a nonzero element of E. Let moreover \(\widetilde{f}_2 \in \mathrm{E}\) non-zero. The map \(h: g \mapsto \langle \widetilde{f}_{1,\mathrm{E}} | \gamma_{\mathrm{G},\overline{\chi}}(g)\widetilde{f}_2 \rangle\) is a matrix coefficient of \(\pi\) . For every \(g \in \mathrm{G}\) , we have \(h(g) = \langle \widetilde{f}_1 | \gamma_{\mathrm{G},\overline{\chi}}(g)\widetilde{f}_2 \rangle\) since \(\widetilde{f}_1 - \widetilde{f}_{1,\mathrm{E}}\) is orthogonal to E. By Lemma 8 of V, p. 418, the function \(h\) belongs to \(\mathcal{L}_{\chi}^{2}(\mathrm{G})\) , so Proposition 6 implies that \(\pi\) is square-integrable modulo the center.

Conversely, suppose that \(\pi\) is square-integrable modulo the center. Let \(x_0 \in \mathrm{E}\) a nonzero vector. By prop. 7 and formula (10), the map \(y \mapsto \check{f}_{x_0,y}\) is an injective G-morphism from \(\pi\) in \(\gamma_{\mathrm{G},\overline{\chi}}\) .

DEFINITION 4. --- Let Z be a closed subgroup of C such that C/Z is compact. Let π be an irreducible unitary representation of G which is square-integrable modulo the center. The unique real number c \textgreater{} 0 satisfying the property of Proposition 7 is called the formal degree of π relative to Z. It is denoted by \(c_{\mathrm{Z}}(\pi)\) .

The formal degree depends on the choice of Haar measure on Z. If the Haar measure \(\beta_{Z}\) on Z is multiplied by a real number t \textgreater{} 0, then the measure \(\nu_{Z} = \mu/\beta_{Z}\) on G/Z is multiplied by \(t^{-1}\) , and the formal degree of \(\pi\) is multiplied by t.

The formal degree is characterized by the following property:

PROPOSITION 8 (Orthogonality relations). --- Let π be an irreducible unitary representation of G in a Hilbert space E that is square-integrable modulo the center. Then one has

\[
c _ {\mathrm{Z}} (\pi) \int_ {\mathrm{G} / \mathrm{Z}} \overline {{\langle x | \pi (g) x ^ {\prime} \rangle}} \langle y | \pi (g) y ^ {\prime} \rangle d \nu_ {\mathrm{Z}} (g) = \overline {{\langle x | y \rangle}} \langle x ^ {\prime} | y ^ {\prime} \rangle
\]

for all \((x,y,x',y')\in \mathrm{E}^4\)

Denote by \(w\) the morphism of Prop. 7. We have

\[
\int_ {\mathrm{G} / \mathrm{Z}} \overline {{\langle x | \pi (g) x ^ {\prime} \rangle}} \langle y | \pi (g) y ^ {\prime} \rangle d \nu_ {\mathrm{Z}} (g) = \langle f _ {x, x ^ {\prime}} | f _ {y, y ^ {\prime}} \rangle
\]

and according to loc. cit., it follows that

\[
\begin{array}{c} \langle f _ {x, x ^ {\prime}} | f _ {y, y ^ {\prime}} \rangle = \frac {1}{c _ {\mathrm{Z}} (\pi)} \langle w (x \otimes x ^ {\prime}) | w (y \otimes y ^ {\prime}) \rangle \\ = \frac {1}{c _ {\mathrm{Z}} (\pi)} \langle x \otimes x ^ {\prime} | y \otimes y ^ {\prime} \rangle = \frac {1}{c _ {\mathrm{Z}} (\pi)} \langle x | y \rangle_ {\overline {{\mathrm{E}}}} \langle x ^ {\prime} | y ^ {\prime} \rangle_ {\mathrm{E}}, \end{array}
\]

whence the result.

In addition to the previous proposition, we also have the following relations for non-isomorphic square-integrable irreducible representations.

PROPOSITION 9. --- Let \(\pi_1\) and \(\pi_2\) non-isomorphic irreducible unitary representations of G in Hilbert spaces E \(_1\) and E \(_2\) . One assumes that \(\pi_1\) and \(\pi_2\) are square-integrable modulo the center and the restrictions to Z of their central characters coincide. Then

\[
\int_ {\mathrm{G} / \mathrm{Z}} \overline {{\langle x \mid \pi_ {1} (g) x ^ {\prime} \rangle}} \langle y \mid \pi_ {2} (g) y ^ {\prime} \rangle d \nu_ {\mathrm{Z}} (g) = 0
\]

for all \((x,x^{\prime},y,y^{\prime})\in \mathrm{E}_{1}^{2}\times \mathrm{E}_{2}^{2}\)

 For \(i=1,2\) , denote \(w_{i}\) the morphism from Proposition 7 for the representation \(\pi_{i}\) . By Lemma 8, \(b\) ) of V, p. 384 and assertion \(b\) ) of Proposition 8 of V, p. 394, the image of \(w_{i}\) is contained in the component \(\pi_{i}\) -isotypic of \(\delta_{\mathrm{G},\chi}\) . By assertion \(a\) ) of loc. cit., the image of \(w_{1}\) is therefore orthogonal to the image of \(w_{2}\) . Consequently, it follows that \(\langle w_{1}(x\otimes x^{\prime})|w_{2}(y\otimes y^{\prime})\rangle=0\) for all \((x,x^{\prime},y,y^{\prime})\in\mathrm{E}_{1}^{2}\times\mathrm{E}_{2}^{2}\) , which is the desired formula.

Remark. --- The orthogonality relations of Props. 8 and 9 generalize those of A, VIII, p. 399 (see also the case where G is compact in § 2 of V, p. 457).

